\documentclass[12pt,a4paper]{report}

% ── Encodage & langue ──────────────────────────────────────────────
\usepackage[utf8]{inputenc}
\usepackage[T1]{fontenc}
\usepackage[french]{babel}

% ── Mise en page ───────────────────────────────────────────────────
\usepackage[top=2.5cm, bottom=2.5cm, left=2.5cm, right=2cm]{geometry}
\usepackage{setspace}
\onehalfspacing

% ── Typographie ────────────────────────────────────────────────────
\usepackage{lmodern}
\usepackage{microtype}

% ── Couleurs & graphiques ──────────────────────────────────────────
\usepackage[dvipsnames,table]{xcolor}
\usepackage{graphicx}
\usepackage{tikz}

\definecolor{accent}{RGB}{99,102,241}
\definecolor{accentdark}{RGB}{67,56,202}
\definecolor{lightgray}{RGB}{248,249,250}
\definecolor{codebg}{RGB}{244,244,244}
\definecolor{bordercolor}{RGB}{220,220,220}

% ── En-têtes & pieds de page ───────────────────────────────────────
\usepackage{fancyhdr}
\pagestyle{fancy}
\fancyhf{}
\fancyhead[L]{\small\textcolor{gray}{Audio Understanding --- Qwen2-Audio-7B-Instruct}}
\fancyhead[R]{\small\textcolor{gray}{Tassnim \& Zineb}}
\fancyfoot[C]{\thepage}
\renewcommand{\headrulewidth}{0.4pt}

% ── Liens hypertexte ───────────────────────────────────────────────
\usepackage[colorlinks=true,linkcolor=accentdark,urlcolor=accentdark,citecolor=accentdark]{hyperref}
\usepackage{url}

% ── Code source ────────────────────────────────────────────────────
\usepackage{listings}
\lstset{
  backgroundcolor=\color{codebg},
  basicstyle=\ttfamily\small,
  breaklines=true,
  breakatwhitespace=true,
  frame=leftline,
  framerule=3pt,
  rulecolor=\color{accent},
  numbers=left,
  numberstyle=\tiny\color{gray},
  numbersep=8pt,
  showstringspaces=false,
  keywordstyle=\color{accentdark}\bfseries,
  stringstyle=\color{OliveGreen},
  commentstyle=\color{gray}\itshape,
  captionpos=b,
  xleftmargin=1.5em,
}
\lstdefinestyle{json}{
  language={},
  morestring=[b]",
  stringstyle=\color{OliveGreen},
}
\lstdefinestyle{bash}{
  language=bash,
  morekeywords={curl, uvicorn, pip, python},
}
\lstdefinestyle{python}{
  language=Python,
  morekeywords={True, False, None, from, import, as, with, yield, async, await},
  stringstyle=\color{OliveGreen},
}

% ── Tableaux ───────────────────────────────────────────────────────
\usepackage{booktabs}
\usepackage{array}
\usepackage{tabularx}
\usepackage{colortbl}

% ── Titres de sections ─────────────────────────────────────────────
\usepackage{titlesec}
\titleformat{\chapter}[block]
  {\normalfont\Large\bfseries\color{accentdark}}
  {\thechapter.}{.6em}{}[\vspace{-0.4em}\rule{\linewidth}{1.5pt}]
\titleformat{\section}
  {\normalfont\large\bfseries\color{accentdark}}
  {\thesection}{.5em}{}
\titleformat{\subsection}
  {\normalfont\normalsize\bfseries}
  {\thesubsection}{.5em}{}

% ── Boîtes colorées ────────────────────────────────────────────────
\usepackage{tcolorbox}
\tcbuselibrary{skins,breakable}
\newtcolorbox{infobox}[1][]{
  enhanced, breakable,
  colback=lightgray, colframe=accent,
  leftrule=4pt, rightrule=0pt, toprule=0pt, bottomrule=0pt,
  arc=2pt, boxsep=4pt,
  fonttitle=\bfseries\small, title=#1,
}

% ── Diagrammes avec verbatim ───────────────────────────────────────
\usepackage{verbatim}
\usepackage{fancyvrb}

% ── Bibliographie ──────────────────────────────────────────────────
\usepackage[backend=biber,style=numeric,sorting=none]{biblatex}
\addbibresource{references.bib}

% ── Utilitaires ───────────────────────────────────────────────────
\usepackage{enumitem}
\setlist[itemize]{label=\textcolor{accent}{\textbullet}, leftmargin=1.5em}
\usepackage{caption}
\captionsetup{font=small,labelfont=bf}

% ═══════════════════════════════════════════════════════════════════
\begin{document}
% ═══════════════════════════════════════════════════════════════════

% ── Page de garde ─────────────────────────────────────────────────
\begin{titlepage}
  \centering
  \vspace*{1.5cm}

  {\Large\bfseries Université Euromed de Fès}\\[0.4em]
  {\large École d'Ingénierie Digitale et d'Intelligence Artificielle}\\[0.2em]
  {\normalsize 2\textsuperscript{ème} Année --- Intelligence Artificielle}

  \vspace{1cm}
  \rule{\linewidth}{1pt}
  \vspace{0.5cm}

  {\small\itshape Projet de Fin du Module}\\[0.3em]
  {\normalsize Architecture Modernes de Réseaux de Neurones}

  \vspace{1.5cm}

  {\Huge\bfseries\color{accentdark} Rapport de Projet}\\[0.6em]
  {\LARGE Audio Understanding}\\[0.4em]
  {\large Compréhension Audio par Intelligence Artificielle}

  \vspace{1cm}
  \begin{tcolorbox}[
    enhanced, colback=white, colframe=accent,
    width=0.65\linewidth, arc=4pt,
    halign=center
  ]
    \centering\ttfamily\large Qwen2-Audio-7B-Instruct
  \end{tcolorbox}

  \vspace{2cm}

  \begin{tabular}{rl}
    \textbf{Étudiantes :}  & Tassnim \& Zineb \\[0.3em]
    \textbf{Encadrant :}   & Pr. Smail Tigani \\[0.3em]
    \textbf{Année :}       & 2025 / 2026 \\
  \end{tabular}

  \vfill
  \rule{\linewidth}{0.5pt}
\end{titlepage}

% ── Table des matières ─────────────────────────────────────────────
\tableofcontents
\newpage

% ═══════════════════════════════════════════════════════════════════
\chapter{Introduction}
% ═══════════════════════════════════════════════════════════════════

\section{Contexte et motivation}

L'intelligence artificielle connaît depuis quelques années une révolution majeure avec
l'émergence des modèles de langage de grande taille (\textit{Large Language Models}, LLMs)
et, plus récemment, des modèles multimodaux capables de traiter simultanément plusieurs
types de données. Parmi ces avancées, la compréhension audio par IA représente une rupture
technologique significative : il ne s'agit plus seulement de transcrire la parole, mais de
\textit{comprendre} le contenu sonore dans toute sa richesse sémantique, émotionnelle et
contextuelle.

Ce projet s'inscrit dans le cadre du module \textit{Architecture Modernes de Réseaux de
Neurones} dispensé à l'École d'Ingénierie Digitale et d'Intelligence Artificielle de
l'Université Euromed de Fès. Il vise à concevoir et déployer une plateforme opérationnelle
de compréhension audio, en s'appuyant sur un modèle open source de pointe :
\textbf{Qwen2-Audio-7B-Instruct}, développé par Alibaba Cloud et disponible sur la
plateforme Hugging Face.

La pertinence de ce sujet est indéniable dans un monde où les interfaces vocales, les
assistants intelligents et les outils d'analyse de réunions se multiplient. Maîtriser la
chaîne complète --- du modèle brut à l'API exploitable --- constitue une compétence
fondamentale pour tout ingénieur en intelligence artificielle.

\section{Objectifs du projet}

Le projet poursuit trois objectifs principaux, alignés sur les critères d'évaluation du
module :

\begin{enumerate}
  \item \textbf{Monter une API opérationnelle} exposant des points de terminaison POST
        capables de recevoir des fichiers audio, de les traiter via le modèle
        Qwen2-Audio-7B-Instruct et de retourner une réponse textuelle structurée.

  \item \textbf{Développer une mise en application concrète} matérialisée par une interface
        web interactive permettant la transcription, la compréhension (Q\&A) et l'analyse
        approfondie de contenus audio.

  \item \textbf{Acquérir une expertise pratique} sur les modèles multimodaux audio-langage
        en comprenant leur architecture, leur intégration et leurs limites.
\end{enumerate}

% ═══════════════════════════════════════════════════════════════════
\chapter{Méthodologie}
% ═══════════════════════════════════════════════════════════════════

\section{Présentation du modèle Qwen2-Audio-7B-Instruct}

\textbf{Qwen2-Audio-7B-Instruct} est un modèle multimodal de compréhension audio développé
par le groupe Qwen d'Alibaba Cloud \cite{qwen2audio}. Il est publié en open source sous
licence Apache 2.0 sur Hugging Face et constitue la version affinée par instruction
(\textit{instruction-tuned}) du modèle de base Qwen2-Audio-7B.

\subsection{Architecture interne}

Le modèle repose sur une architecture hybride associant deux composants distincts :

\begin{itemize}
  \item \textbf{Un encodeur audio} de type Whisper Large V2 \cite{whisper}, responsable de
        l'extraction des représentations acoustiques à partir des formes d'onde brutes.
        Cet encodeur traite le signal audio échantillonné à 16\,000 Hz et produit une
        séquence de vecteurs de caractéristiques denses.

  \item \textbf{Un décodeur de langage Qwen2-7B} \cite{qwen2}, un modèle de langage
        autorégressif à 7 milliards de paramètres, qui reçoit en entrée la concaténation
        des représentations audio et des tokens textuels, et génère la réponse token par
        token.
\end{itemize}

\begin{figure}[h]
\centering
\begin{Verbatim}[frame=single, rulecolor=\color{accent}, fontsize=\small]
Audio (WAV 16 kHz)
       |
       v
+------------------------------+
|   Encodeur Audio             |
|   (Whisper Large V2)         |
|   --> Feature Extraction     |
|   --> Audio Embeddings       |
+-------------+----------------+
              |   Audio tokens
              v
+------------------------------+      Prompt texte
|   Fusion Multimodale         | <--- (question/instruction)
|   Concatenation des tokens   |
+-------------+----------------+
              |
              v
+------------------------------+
|   Decodeur LLM Qwen2-7B      |
|   (7B parametres)            |
|   Generation autoregressive  |
+-------------+----------------+
              |
              v
       Reponse textuelle
\end{Verbatim}
\caption{Architecture du modèle Qwen2-Audio-7B-Instruct}
\end{figure}

\subsection{Capacités du modèle}

Qwen2-Audio-7B-Instruct excelle dans plusieurs tâches audio-langage : la reconnaissance
automatique de la parole (ASR) en plusieurs langues, la réponse à des questions libres
portant sur le contenu sonore, l'analyse émotionnelle et tonale, ainsi que la description
d'événements sonores non-verbaux (musique, bruits ambiants, etc.).

\section{Architecture du système}

Le système développé s'articule autour de deux couches principales : un backend Python
assurant l'inférence du modèle et l'exposition des endpoints, et un frontend HTML/CSS/JS
offrant l'interface utilisateur.

\begin{figure}[h]
\centering
\begin{Verbatim}[frame=single, rulecolor=\color{accent}, fontsize=\small]
+----------------------------------------------+
|               Navigateur Web                 |
|  Interface SPA -- 3 onglets :                |
|  [ Transcrire ] [ Comprendre ] [ Analyser ]  |
|  Upload fichier  /  Enregistrement micro     |
+--------------------+-------------------------+
                     | HTTP POST multipart/form-data
                     v
+----------------------------------------------+
|         Backend FastAPI (port 8000)          |
|                                              |
|  POST /api/v1/audio/transcribe               |
|  POST /api/v1/audio/understand               |
|  POST /api/v1/audio/analyze                  |
|  GET  /api/v1/health                         |
+--------------------+-------------------------+
                     | Inference locale
                     v
+----------------------------------------------+
|     Qwen2-Audio-7B-Instruct (HuggingFace)    |
|     CUDA float16  /  CPU float32             |
+----------------------------------------------+
\end{Verbatim}
\caption{Architecture globale du système}
\end{figure}

\section{Déploiement sur Google Colab avec quantification 4 bits}

Faute de GPU local suffisant, le système est déployé sur \textbf{Google Colab} en
exploitant un GPU \textbf{NVIDIA T4} (15~Go de VRAM), le GPU gratuit le plus courant
sur la plateforme.

\subsection{Pourquoi la quantification est indispensable}

Le modèle Qwen2-Audio-7B-Instruct comporte 7 milliards de paramètres. En précision
\texttt{float16} (2 octets par paramètre), le chargement seul occupe environ
$7 \times 10^9 \times 2 = 14~\mathrm{Go}$, auxquels s'ajoutent les activations
intermédiaires et les buffers du système CUDA, portant le besoin réel à
\textbf{$\approx$16~Go de VRAM}. Le T4 ne dispose que de 15~Go~: le chargement en
float16 est donc impossible sans débordement mémoire.

La \textbf{quantification en 4 bits NF4} (\textit{NormalFloat-4}), fournie par la
bibliothèque \texttt{bitsandbytes}, résout ce problème en ne stockant chaque paramètre
que sur un demi-octet, ramenant l'empreinte mémoire à \textbf{$\approx$4--5~Go} :

\begin{table}[h]
\centering
\begin{tabular}{lccc}
\toprule
\textbf{Précision} & \textbf{Octets/paramètre} & \textbf{VRAM requise} & \textbf{GPU Colab compatible} \\
\midrule
float16 (FP16)   & 2,0 & $\approx$16 Go & A100 (40 Go) \\
8 bits (INT8)    & 1,0 & $\approx$8 Go  & A100, V100 \\
\textbf{4 bits NF4} & \textbf{0,5} & \textbf{$\approx$5 Go} & \textbf{T4, tout GPU $\geq$ 6 Go} \\
\bottomrule
\end{tabular}
\caption{Empreinte mémoire du modèle selon la précision}
\end{table}

NF4 est un format spécialement conçu pour les distributions gaussiennes typiques des
poids des réseaux de neurones. Combiné à la \textit{double quantification} (qui
quantifie elle-même les constantes de quantification) et à un \textit{compute dtype}
en \texttt{float16} (les multiplications matricielles restent en demi-précision), la
dégradation de qualité reste minimale pour les tâches de compréhension audio.

Dans le code (\texttt{backend/app/model.py}), la quantification est activée via la
variable d'environnement \texttt{LOAD\_IN\_4BIT=true} :

\begin{lstlisting}[style=python, caption={Configuration BitsAndBytesConfig pour la quantification NF4}]
from transformers import BitsAndBytesConfig

quant_config = BitsAndBytesConfig(
    load_in_4bit=True,
    bnb_4bit_quant_type="nf4",          # NormalFloat-4
    bnb_4bit_compute_dtype=torch.float16,# calculs en fp16
    bnb_4bit_use_double_quant=True,      # double quantification
)
model = Qwen2AudioForConditionalGeneration.from_pretrained(
    MODEL_NAME,
    quantization_config=quant_config,
    device_map="auto",
)
\end{lstlisting}

\subsection{Procédure de déploiement}

La procédure est automatisée par le notebook \texttt{Qwen2\_Audio\_Colab.ipynb} :

\begin{enumerate}
  \item Sélection du GPU T4 dans Colab (\textit{Runtime → Change runtime type}).
  \item Clonage du dépôt GitHub du projet dans l'environnement Colab.
  \item Installation des dépendances (\texttt{transformers}, \texttt{accelerate},
        \texttt{bitsandbytes}, \texttt{librosa}, \texttt{ffmpeg}, FastAPI) sans
        réinstaller PyTorch, déjà fourni avec CUDA par Colab.
  \item Activation de la variable \texttt{LOAD\_IN\_4BIT=true} avant le lancement.
  \item Lancement du serveur \texttt{uvicorn} sur le port 8000.
  \item Ouverture d'un tunnel public via \texttt{cloudflared}, qui expose
        l'application sur une URL temporaire \texttt{*.trycloudflare.com}.
\end{enumerate}

Comme le serveur FastAPI sert à la fois l'interface web et l'API sur la même origine,
l'ouverture de cette unique URL publique suffit à utiliser la plateforme complète
depuis n'importe quel navigateur. Le premier démarrage télécharge environ 4--5~Go de
poids quantifiés (contre 16~Go en float16) ; les exécutions suivantes utilisent le
cache Hugging Face local.

\section{API opérationnelle (FastAPI)}

Le backend est développé avec le framework \textbf{FastAPI} \cite{fastapi} (Python),
reconnu pour ses performances élevées, sa validation automatique des données via Pydantic,
et sa génération automatique de documentation interactive (Swagger UI).

\subsection{Endpoints implémentés}

\begin{table}[h]
\centering
\rowcolors{2}{lightgray}{white}
\begin{tabularx}{\textwidth}{llXX}
\toprule
\textbf{Méthode} & \textbf{Endpoint} & \textbf{Entrée} & \textbf{Sortie} \\
\midrule
POST & \texttt{/api/v1/audio/transcribe} & Fichier audio & Transcription complète \\
POST & \texttt{/api/v1/audio/understand} & Fichier audio + question & Réponse à la question \\
POST & \texttt{/api/v1/audio/analyze}    & Fichier audio & Transcription + résumé + sentiment \\
GET  & \texttt{/api/v1/health}           & ---           & État du modèle et du serveur \\
\bottomrule
\end{tabularx}
\caption{Endpoints de l'API}
\end{table}

\subsection{Pipeline de traitement}

Chaque requête suit le pipeline suivant :

\begin{enumerate}
  \item Réception et validation du fichier audio (formats : WAV, MP3, OGG, FLAC, WebM,
        M4A ; taille maximale : 50 Mo).
  \item Décodage et rééchantillonnage à 16\,000 Hz via \texttt{soundfile} et
        \texttt{librosa} \cite{librosa}.
  \item Construction de la conversation au format \textit{chat template} Qwen2-Audio.
  \item Appel au modèle via \texttt{model.generate()}.
  \item Sérialisation de la réponse en JSON.
\end{enumerate}

\subsection{Exemple de requête et réponse}

\begin{lstlisting}[style=bash, caption={Requête cURL vers /understand}]
curl -X POST http://localhost:8000/api/v1/audio/understand \
  -F "audio=@reunion.wav" \
  -F "question=Quels sont les points principaux abordes ?"
\end{lstlisting}

\begin{lstlisting}[style=json, caption={Réponse JSON de l'endpoint /understand}]
{
  "question": "Quels sont les points principaux abordes ?",
  "answer": "La reunion porte principalement sur le calendrier du
             projet, la repartition des taches et le budget alloue
             pour le trimestre.",
  "filename": "reunion.wav"
}
\end{lstlisting}

\section{Interface web}

L'interface utilisateur est une application monopage (\textit{Single Page Application})
développée en HTML5, CSS3 et JavaScript vanilla, servie directement par le backend FastAPI.
Elle s'organise en trois onglets :

\begin{itemize}
  \item \textbf{Transcrire} --- l'utilisateur soumet un fichier audio ou enregistre sa
        voix directement via le microphone ; le système retourne la transcription intégrale.

  \item \textbf{Comprendre} --- l'utilisateur joint un fichier audio et formule une
        question en langage naturel ; le modèle répond en s'appuyant sur le contenu de
        l'audio.

  \item \textbf{Analyser} --- le système produit automatiquement un rapport complet :
        transcription, résumé des points clés et analyse du ton/sentiment exprimé.
\end{itemize}

L'enregistrement audio est réalisé via l'API Web \texttt{MediaRecorder} du navigateur,
permettant une interaction directe sans logiciel tiers. Le fichier produit est transmis
en \texttt{multipart/form-data} vers le backend.

% ═══════════════════════════════════════════════════════════════════
\chapter{Résultats}
% ═══════════════════════════════════════════════════════════════════

\section{Fonctionnalités implémentées}

L'ensemble des fonctionnalités prévues a été implémenté avec succès. Le tableau
ci-dessous récapitule les livrables du projet.

\begin{table}[h]
\centering
\rowcolors{2}{lightgray}{white}
\begin{tabularx}{\textwidth}{llX}
\toprule
\textbf{Livrable} & \textbf{Statut} & \textbf{Détails} \\
\midrule
API POST /transcribe    & Implémenté & Multilingue, WAV/MP3/OGG/FLAC/WebM/M4A \\
API POST /understand    & Implémenté & Question libre en langage naturel \\
API POST /analyze       & Implémenté & Transcription + résumé + sentiment \\
Interface web           & Implémenté & 3 onglets, upload + enregistrement micro \\
Documentation Swagger   & Auto-générée & Accessible à \texttt{/api/docs} \\
Validation des entrées  & Implémenté & Format, taille, contenu vide \\
Mode développement      & Implémenté & Variable \texttt{USE\_MOCK} sans GPU \\
\bottomrule
\end{tabularx}
\caption{Récapitulatif des livrables}
\end{table}

\section{Exemples de sorties}

Les tests ont été réalisés avec des fichiers audio variés (parole en français et en
anglais, enregistrements de réunions, extraits musicaux).

\begin{infobox}[Transcription --- fichier audio de 30 secondes en français]
\textit{Sortie :} « Bonjour à tous. Aujourd'hui nous allons discuter des résultats du
dernier trimestre. Nos indicateurs montrent une progression de 15\,\% par rapport à
l'année précédente, ce qui dépasse nos objectifs initiaux... »
\end{infobox}

\begin{infobox}[Compréhension (Q\&A) --- Question : « Quelle est l'émotion du locuteur ? »]
\textit{Réponse :} « Le locuteur exprime une satisfaction marquée et un enthousiasme
modéré. Son ton est professionnel et confiant, avec des inflexions positives lorsqu'il
mentionne les résultats dépassant les objectifs. »
\end{infobox}

\begin{infobox}[Analyse complète --- résumé et sentiment]
\textit{Résumé :} « L'audio présente un discours de bilan trimestriel. Points clés :
progression de 15\,\%, dépassement des objectifs, annonce d'un nouveau plan. »\\[0.3em]
\textit{Sentiment :} « Ton positif et professionnel. Sentiment dominant : satisfaction et
optimisme. »
\end{infobox}

\section{Difficultés rencontrées}

Le développement du projet a présenté plusieurs défis techniques :

\begin{itemize}
  \item \textbf{Gestion des formats audio hétérogènes.} Les navigateurs web produisent des
        enregistrements au format WebM/Opus, non nativement supporté par
        \texttt{soundfile}. La solution adoptée consiste à utiliser \texttt{librosa} comme
        \textit{fallback}, qui s'appuie sur \texttt{ffmpeg} pour décoder les formats
        exotiques.

  \item \textbf{Contrainte mémoire GPU.} Le modèle Qwen2-Audio-7B-Instruct requiert
        environ 16~Go de VRAM en précision \texttt{float16}, ce qui dépasse la capacité
        du GPU T4 disponible sur Google Colab (15~Go). La solution retenue est la
        \textbf{quantification NF4 4 bits} via \texttt{bitsandbytes}, qui ramène
        l'empreinte mémoire à environ 4--5~Go tout en conservant une qualité de
        transcription et de compréhension satisfaisante. La variable d'environnement
        \texttt{LOAD\_IN\_4BIT=true} active cette configuration au démarrage.

  \item \textbf{Audio non transmis au modèle.} Le défi technique le plus subtil a été
        un bug où le modèle répondait comme un simple modèle de langage textuel
        (« je ne peux pas accéder à l'audio »), sans aucune erreur apparente. La cause :
        l'argument du processeur transportant l'audio a été \textbf{renommé} d'une
        version de \texttt{transformers} à l'autre (\texttt{audios} $\rightarrow$
        \texttt{audio}). Lorsqu'on passe le mauvais nom, l'argument est silencieusement
        ignoré (absorbé par \texttt{**kwargs}) et l'audio disparaît des entrées du modèle.
        \textbf{Solution :} essayer successivement les deux noms et ne conserver que
        l'appel produisant effectivement des caractéristiques audio
        (\texttt{input\_features}), échouer explicitement sinon, et aligner le type des
        tenseurs audio (float16) sur celui du modèle.

  \item \textbf{Gestion du chat template multimodal.} La méthode
        \texttt{apply\_chat\_template} du processeur Qwen2-Audio nécessite une structure
        de conversation précise mélangeant tokens audio et tokens textuels, dont la
        compréhension a demandé une lecture approfondie du code source du modèle.

  \item \textbf{Temps d'initialisation.} Le chargement du modèle depuis Hugging Face
        représente environ 16 Go de téléchargement à la première exécution. Le modèle est
        ensuite mis en cache localement pour les exécutions suivantes.
\end{itemize}

% ═══════════════════════════════════════════════════════════════════
\chapter{Conclusion}
% ═══════════════════════════════════════════════════════════════════

\section{Bilan du projet}

Ce projet nous a permis de maîtriser la chaîne de valeur complète d'un système d'IA
audio en production. Nous sommes parties de la compréhension théorique de l'architecture
Qwen2-Audio --- encodeur Whisper couplé à un décodeur Qwen2-7B --- pour aboutir à une
application web pleinement opérationnelle, exposant trois endpoints REST bien définis et
une interface utilisateur soignée.

Le projet illustre concrètement comment des modèles open source de grande taille peuvent
être déployés dans des scénarios réels sans recourir à des APIs propriétaires tierces.
L'approche adoptée --- inférence locale, API FastAPI, frontend SPA --- constitue une
architecture reproductible et extensible pour de nombreux cas d'usage industriels.

Sur le plan pédagogique, ce projet a renforcé notre compréhension des modèles
multimodaux, de la gestion des flux audio, de la conception d'APIs REST et du
développement d'interfaces web modernes. Il nous a également sensibilisées aux
contraintes pratiques du déploiement de LLMs : consommation mémoire, temps de
chargement et gestion des formats de données hétérogènes.

\section{Pistes d'amélioration}

Plusieurs améliorations pourraient enrichir ce projet dans une version future :

\begin{itemize}
  \item \textbf{Streaming des réponses.} Implémenter les \textit{Server-Sent Events} (SSE)
        pour afficher la réponse du modèle token par token, améliorant ainsi la réactivité
        perçue par l'utilisateur.

  \item \textbf{Conversation multi-tours.} Permettre plusieurs échanges successifs portant
        sur le même fichier audio, en maintenant un contexte de dialogue.

  \item \textbf{Détection de langue automatique.} Identifier automatiquement la langue du
        locuteur et adapter la réponse en conséquence.

  \item \textbf{Support des audios longs.} Implémenter un mécanisme de découpage
        (\textit{chunking}) pour traiter des enregistrements de plusieurs heures (réunions,
        conférences, podcasts).

  \item \textbf{Quantification avancée.} Au-delà de la quantification 4/8 bits
        (\texttt{bitsandbytes}) déjà prise en charge, explorer des méthodes hors ligne
        telles que GPTQ ou AWQ pour réduire davantage l'empreinte mémoire et accélérer
        l'inférence sur des GPUs de faible capacité.

  \item \textbf{Évaluation formelle (WER).} Mesurer le \textit{Word Error Rate} sur un
        corpus francophone de référence pour quantifier objectivement la précision du
        modèle.
\end{itemize}

% ── Bibliographie ──────────────────────────────────────────────────
\printbibliography[heading=bibintoc, title={Références bibliographiques}]

\end{document}
